\section{Budget mechanism details}\label{app:budget}

\subsection{What was wrong with the time bank}
The time bank never gave a task less than its base share, and every finished task donated its unused seconds to all running and queued tasks. That made the outcome depend on queue order: in a simulation of the Kaggle run, early tasks got about 1.1x the base and late ones about 3x. The harness's ``easy first'' order (by grid size) put the historically hardest tasks first (expected solve rate 0.28 against 0.53 for the last quarter). Replaying the exact time-bank code on 240 tasks showed the run ending 0.5--2.4 hours early while 41--49 hard first-wave tasks were cut at 100--115k tokens, against the 127--137k they needed; a value model put the cost of those cuts at no more than 7--12 points per 240 tasks (about 3--5 leaderboard points). On our machines, multi-wave runs under the time bank scored 4--6 below single-wave runs of the same harness, which is the same effect.

\subsection{Design of budget v2}
\begin{itemize}
\item Base 86k generated tokens per task; unused tokens go to a shared pool.
\item A task is flagged hard when its first answer carries a P0/P1 doubt or nothing passes training by 60k tokens. In practice 52--55 of 120 tasks were flagged (10--13 for a doubt, 42 for no pass), and flagged tasks used 111--115k on average.
\item Flagged tasks get a floor of 160k from the pool, extended lazily to 192k when they come within 10k of their allotment. The values came from the hard-task budget curve (Appendix~\ref{app:resource}).
\item A tail queue runs second chains capped at 160k (at least 100k needed to start): fresh conversations that stop at their first train-passing program and replace attempt 2. Order: doubted tasks with long first chains first; tasks without any train-passing first attempt last.
\item On Kaggle the deadline of every task is the global end of the run, tokens set the budgets, and leftover wall time near the end also funds chains.
\end{itemize}

\subsection{Results}
\begin{table}[H]\centering\small
\begin{tabular}{@{}lrrr@{}}
\toprule
Run (multi-wave, 120 tasks, 24 lanes, 86k per task) & augD1 & augD2 & Tokens \\
\midrule
Time bank, run 1 & 94.17 & 91.33 & 9.24M / 8.81M \\
Time bank, run 2 & 91.00 & 93.17 & -- \\
Budget v2 on the v32 harness (primary only) & 98.83 (97.83) & 97.17 (94.17) & 9.06M / 8.96M \\
Budget v2, exact v33 build & 99.50 & 98.67 & -- \\
Budget v2, v33 build, repeat & 93.50 & 94.83 & -- \\
\bottomrule
\end{tabular}
\caption{Budget v2 against the time bank. Four-run means: 98.54 against 92.42 (+6.1). Including the two later repeats of the v33 build, the budget-v2 mean over six runs is 97.1. On the leaderboard v33 scored 73.89 against the v32 mean of 70.14.}
\end{table}

\subsection{Tail chains}
\begin{itemize}
\item Value depends on whom they go to. In a selection study of 175 tail chains, tasks with no train-passing first attempt got 0 rescues from 8 chains of 120k each; tasks with a doubted attempt 1 and a primary of at least 100k gave the best return, 5.0 points from 41 chains (1.33 per million tokens).
\item Length matters: a chain cut at 60k kept 18--26\% of a full chain's value, 85--90\% at 120k. On 45 hard tasks a 64k second chain was worth about nothing, a 128k chain 3.3 points per run.
\item With a realistic pool (about 11k spare tokens per task on Kaggle), the order decides whether tail chains pay: with the old order the net was about 0, hard-first made it positive.
\item Merge rule: replacing attempt 2 scored 31.5 on 57 tail task-runs, against 29.0 for filling only an empty slot, 28.0 per test input and 26.0 for keeping the primary. Weighted voting lost on held-out runs ($-1.0$ per run).
\end{itemize}
